The goal of anomaly detection is to identify deviations from normal behavior. Using labeled anomalies to train or calibrate a detector can bias it toward known anomaly patterns, weakening robustness to unforeseen, ``zero-day'' deviations \cite{ruff_review_2021}. This concern, together with the scarcity of anomaly labels \cite{chandola_review}, motivates learning normal behavior directly and scoring departures from it. Common approaches include probabilistic models that derive scores from observation likelihoods under a reference distribution \cite{probabilistic_ad}; distance- and density-based methods that quantify separation from neighbors \cite{knn_ad} or relative local sparsity \cite{lof_ad}; and reconstruction and forecasting models that derive scores from discrepancies between their outputs and the observed values \cite{dl_tsad_review}.

Deep learning models can capture nonlinear relationships across variables and time, producing anomaly scores that reflect departures from complex patterns of normal behavior \cite{dl_tsad_review}. Recurrent networks and transformers differ in architecture, but both can provide a score stream to a downstream decision rule. The scoring model is therefore treated here as an upstream component rather than the object of the proposed calibration procedure.

This distinction also matters when interpreting reported performance. Threshold-independent metrics such as the area under the ROC curve assess score ranking but do not specify an operational rule for generating anomaly labels. Conversely, an $F_1$ score obtained by maximizing performance over thresholds on the labeled test set evaluates an operating point selected using information unavailable under normal-only calibration \cite{xu_bostrom_2026}. Comparing decision rules therefore requires keeping the score stream fixed and distinguishing evaluation with anomaly labels from calibration without them.
